\section{Introduction}
\label{sec:intro}

%-------------------------------------------------------------------------

\subsection{Animal Localization and Preprocessing}
In this study, subject localization and extraction are performed using a specialized two-stage pipeline designed to isolate the primary animal before fine-grained classification. 
To comply with the strict inference time constraints of the project averaging under 5 seconds per image, while maintaining high spatial accuracy. 
We use an off-the-shelf YOLOv6 object detector (Li et al., 2022a). Specifically, the lightweight YOLOv6-S architecture is used. The YOLOv6 framework introduces significant advancements in 
single-stage detection, utilizing hardware-aware neural network design and reparameterization techniques to maximize throughput without sacrificing performance (Li et al., 2022b). 
Furthermore, it incorporates an adaptive selfdistillation strategy during training, which dynamically adjusts knowledge transfer, rendering it highly robust for real-time edge detection tasks 
(Zhang et al., 2020).To ensure robustness against background interference, our pipeline processes raw RGB images directly, leveraging the detector's native Non-Maximum Suppression (NMS) capabilities. 
At inference, the network outputs bounding box coordinates. Our algorithm explicitly filters these predictions for the target species (cats and dogs, corresponding to COCO IDs 15 and 16). 
In scenarios where multiple target animals are present within a single frame, the bounding box with the largest spatial area is identified as the main subject, directly sticking to the dataset's 
ground-truth definition. A critical challenge in fine-grained classification is contextual confounding, where secondary animals in the background cause false predictions. To mitigate this, we 
implement a zero-out masking technique: all bounding boxes that do not correspond to the primary detections within the frame are explicitly blacked out (pixel values set to zero). 
The primary animal is subsequently cropped and bilinearly resized to a standardized dimension of $224 \times 224$ pixels, meeting the input requirements for subsequent fine-grained 
classification networks. Images yielding no valid detections indicating the presence of confounders or the absence of target species are routed to a fallback mechanism. 
This mechanism passes a resized version of the original frame forward while assigning a 'neither' label, forcing the pipeline to reliably return the required $-1$ reject class.

